\documentclass[11pt]{article}
\usepackage[a4paper,margin=27mm]{geometry}
\usepackage{amsmath,amssymb,amsthm}
\usepackage[hidelinks]{hyperref}
\setlength{\emergencystretch}{3em}
\newtheorem{proposition}{Proposition}
\title{A counterexample to the indexed interlacing claim\\in TLMC conjecture 00000007681}
\author{Local AI-assisted verification package}
\date{21 September 2026}
\begin{document}
\maketitle

\begin{abstract}
The first assertion of TLMC conjecture 00000007681 is false for the
standard continuous $q$-Hermite polynomials. At $q=3/4$, $n=2$, and $k=1$,
the stated inequality requires $0<-1/4$. All indices lie in their stated
root lists. We derive the polynomials from their recurrence and give their
complete, strictly increasing lists of real roots. The accompanying Lean
file formalizes this bridge as well as the contradiction.
\end{abstract}

\section{Statement and interpretation}
The frozen source is conjecture 00000007681 at organizer commit\\
\texttt{95acb520ec5607c826b8a997b1ef2fc82d6f7c57} \cite{tlmc}.
Both language versions give the recurrence
\[
 2xH_n(x\mid q)=H_{n+1}(x\mid q)+(1-q^n)H_{n-1}(x\mid q)
\]
and order the zeros as $x_{n,1}<\cdots<x_{n,n}$.
The first asserted property is
\begin{equation}\label{eq:claim}
 x_{n+1,2k}<x_{n,k}<x_{n+1,2k+1}\qquad(0<q<1).
\end{equation}
We evaluate this only when all three indices exist.

The source does not list initial values. We use the standard continuous
$q$-Hermite normalization $H_0=1$, $H_1=2x$. These initial values follow
from the defining finite sum in NIST DLMF equation 18.28.16 \cite{dlmf}:
for $n=0$ the sum is $1$, and for $n=1$ it is
$e^{i\theta}+e^{-i\theta}=2\cos\theta$.
The paper's use of this published normalization is explicit; the Lean file
starts from these initial polynomials and the displayed recurrence.

\section{Exact counterexample}
\begin{proposition}
Equation \eqref{eq:claim} is false for $q=3/4$, $n=2$, $k=1$.
\end{proposition}
\begin{proof}
The recurrence and initial values give
\[
 H_2(x\mid q)=4x^2-(1-q),\qquad
 H_3(x\mid q)=8x^3-2x(2-q-q^2).
\]
At $q=3/4$ these become
\[
 H_2(x\mid 3/4)=4x^2-\frac14
    =4\left(x+\frac14\right)\left(x-\frac14\right),
\]
\[
 H_3(x\mid 3/4)=x\left(8x^2-\frac{11}{8}\right)
    =8x\left(x+\frac{\sqrt{11}}8\right)
          \left(x-\frac{\sqrt{11}}8\right).
\]
Their exhaustive increasing lists of real zeros are therefore
\[
 (x_{2,1},x_{2,2})=\left(-\frac14,\frac14\right),
\qquad
 (x_{3,1},x_{3,2},x_{3,3})=
   \left(-\frac{\sqrt{11}}8,0,\frac{\sqrt{11}}8\right).
\]
The parameter satisfies $0<3/4<1$. The choice $n=2$, $k=1$ uses the
first zero of $H_2$ and the second and third zeros of $H_3$, so no
out-of-range index is involved. The left-hand inequality in
\eqref{eq:claim} is now $x_{3,2}<x_{2,1}$, or $0<-1/4$, a contradiction.
\end{proof}

\section{Formal coverage and limits}
\begingroup\raggedright
The accompanying Lean source is
\texttt{Results/Counterexample7681.lean}. It defines the polynomial family
over $\mathbb R$, verifies the recurrence and the displayed evaluations,
characterizes all real roots by if-and-only-if statements, and proves that
the two root tables are strictly increasing and exhaustive. The predicate
\texttt{ConjecturedInterlacing} translates \eqref{eq:claim} using zero-based
\texttt{Fin} indices: source index $k$ corresponds to the value $k-1$.
The theorem \texttt{conjectured\_interlacing\_false} negates that predicate.
\par\endgroup

The source also asserts a monotonicity and limiting property of root
differences. We do not formalize that second assertion independently.
The first assertion is necessary for the stated conjunction, so its
counterexample refutes the conjunction regardless of the second
assertion's precise interpretation.

This result diagnoses an indexing defect in the generated statement. It
is not a new general theorem about orthogonal polynomials and does not
refute the correctly indexed ordinary interlacing theorem. A bounded
check of all 100 organizer pull-request titles and bodies, captured at
2026-09-21 01:13:39 UTC+08, found no matching conjecture number. This is
not a claim of global priority. Build, axiom-audit, and replay results are
recorded in the accompanying verification logs; organizer acceptance is
not implied by local checks.

\begin{thebibliography}{9}
\raggedright
\bibitem{tlmc}
The Last Math Competition, conjecture 00000007681,
frozen at commit \texttt{95acb520ec5607c826b8a997b1ef2fc82d6f7c57}.
\url{https://github.com/The-Last-Math-Competition/The-Last-Math-Competition/blob/95acb520ec5607c826b8a997b1ef2fc82d6f7c57/conjectures/00000007681.md}.
\bibitem{dlmf}
NIST Digital Library of Mathematical Functions,
Section 18.28(vi), equation 18.28.16, Continuous $q$-Hermite Polynomials.
\url{https://dlmf.nist.gov/18.28.E16}. Accessed 21 September 2026.
\end{thebibliography}
\end{document}
